{\small\textbf{Credit and housing (Zheng, 1:15) [6:35–7:50]}}\par\smallskip
In the financial system, the picture was split.
\begin{itemize}\setlength\itemsep{0pt}
\item \textbf{Housing was hot.} In early 2022 nominal house prices rose 19 percent year-on-year, and real prices were 16 percent above their 2007 peak. DNB itself called the housing market "overheated".
\item But look at what happened next: real prices fell by about nine percent within a year, as mortgage rates rose.
\item \textbf{Credit was cold.} Household debt is high by international standards, about 107 percent of GDP against 58 percent in the euro area, but it was \textbf{falling}.
\item \textcolor{navy}{\textbf{[def]}} The debt-service ratio, the share of income that goes on interest and repayments, was at its \textbf{lowest since 2005}.
\end{itemize}
\par\smallskip
So if there was a boom, it was in house prices, not in credit.
